\section{Introduction}
\label{sec:intro}

Omni-modal large language models take text, images, audio, and video in one prompt and answer in free text~\cite{gpt-5.2, gemini3, Xu2025Qwen3OmniTR, yu2025minicpm, xu2025qwen2, ye2026omnivinci}.
When such a model answers a question about a forty-second clip or a fifteen-quote document pool, a user who wants to check the answer has to go through the whole input again. A useful system would point to the evidence: these three seconds of audio, this chart, this passage.
We study this problem as \emph{source attribution}: for each generated sentence, find the parts of the input that the model relied on when it produced that sentence.

Three kinds of methods exist, and none of them covers this setting.
\emph{Perturbation methods} such as leave-one-out and ContextCite~\cite{cohen-wang2024contextcite} hide parts of the context and measure how the output changes. They are faithful by construction but need tens of extra forward passes per query, and a video or a long audio clip makes each pass expensive.
\emph{Learned methods} such as AT2~\cite{cohenwang2025at2} and vStream~\cite{kang2026vstream} fit an attribution model on perturbation data, which has to be collected again for every model.
Both families were built for text contexts or for single images, and they have not been tested on continuous audio and video.
\emph{Self-citation} asks the model to state its own evidence. Recent work shows that stated audio and video citations are often hallucinated~\cite{wan2026murgat}, and a citation says what the model claims, not what it used.

Reading attention during generation is the natural cheap alternative: it needs no extra pass and no training. Attention has been used this way for text contexts~\cite{ding2025attention, wang2026attntrace}, and a natural design for mixed inputs traces every generated token to the single input position with the highest attention and then votes over tokens.
We find that this design does not hold up once the input is not text. With hard per-token assignment, the ranking of video seconds agrees with human evidence no better than chance, the question text wins almost every assignment when it competes with video tokens, and a rule tuned on one model loses most of its effect on another (\Cref{sec:design}).
There is also an evaluation gap. Attribution on audio and video has been scored by agreement with reference labels, which shows what a person would point to, not what the model depended on. A causal test needs a way to hide a time interval without leaving a trace of it, and a control for how much was hidden.

We introduce \textbf{OmniTrace}, a training-free attribution method that runs inside the generation pass, together with a causal protocol to test it on every modality.
OmniTrace splits the input into \emph{source units} by structure the input already has: one-second units for audio and video, and one unit per text passage or image.
For each generated token it sums the attention of the late decoder layers over the positions of each unit. Summing these \emph{attention masses} over the tokens of a sentence gives one score per unit, and a unit is labelled as evidence when its score stands out from the other units of its modality.
There is no per-token argmax, no token weighting, no learned component, and no per-modality tuning: the layer band is a fixed fraction of the model's depth and one threshold is used throughout.
The same pass that generates the response returns its attribution, as source ids for text and images and as time intervals for audio and video.

To test whether these units are the ones the model used, we replay the model's own response with units hidden and measure the loss in log-probability.
Hidden units are replaced at the input by placeholders of the same length that carry no content (black frames, silence, a black image, blank tokens), and every result is compared with hiding the same number of random units.
Leave-one-out and ContextCite are run on the same units with the same placeholders, so no baseline is ``not applicable'' on audio or video.

On Qwen2.5-Omni~\cite{xu2025qwen2} and MiniCPM-o~4.5~\cite{yu2025minicpm}, with video (NExT-GQA), audio (MMAU), and image--text documents (MMDocRAG), we find:
\begin{itemize}
    \item \textbf{One pass is enough for most of the faithfulness.} The units OmniTrace selects matter more than random units in all settings, for whole responses and for single sentences. OmniTrace has about twice the effect of embedding similarity and 77--98\% of the effect of leave-one-out and ContextCite, at none of their 16 to 41 extra forward passes.
    \item \textbf{What the model cites is not what it uses.} When asked to cite, the model cites nothing in 18--69\% of responses. On audio and image--text inputs, the cited units matter 1.7 to 3.1 times less than the units OmniTrace selects on the same responses.
    \item \textbf{The simple rule transfers; hard tracing does not.} The same setting works on a second model family without retuning, while per-token argmax tracing loses almost all of its effect on audio there.
    \item \textbf{Attention has measurable limits.} Summed attention favors large units, effects on video are small, and on one model human-marked evidence matters more than the units OmniTrace picks. We quantify each of these.
\end{itemize}

In short, we contribute a single-pass, training-free attribution method for text, image, audio, and video sources, a causal evaluation protocol for continuous inputs, and an analysis of stated versus used evidence, unit-size bias, and modality reliance.
